\section{Operator mother space as a dictionary: resolvents, determinants, and auditable updates}
\label{app:operator_mother_space_dictionary}

\paragraph{Scope and status.}
\InterfaceTag This section introduces a compact dictionary layer for the dynamics narrative: it packages several recurring generating-function viewpoints (Abel-first weights, pole barriers, prime-cycle products, and CAP selection on finite families) into a single operator-theoretic bookkeeping object.
\AuditTag It is reader-facing and does not serve as a premise in theorem-level folding proofs.
For the global cross-module vocabulary used for navigation, see Section~\ref{sec:unified_spine}.
For the shared determinant/trace/pressure bookkeeping identities and invariance contracts used by this dictionary, see Section~\ref{sec:rigidity_bridge_spine}.
For the detailed operator notes, see Appendix~\ref{app:operator_mother_space}.
For the Abel finite-part template and the pole-barrier rigidity template used throughout the paper, see Appendix~\ref{app:abel_finite_part_notes} and Appendix~\ref{app:trace_pole_barrier_template}.

\subsection{Minimal objects: a source operator, readout kernels, and a determinant}
\label{subsec:operator_mother_space_dictionary_minimal}

\paragraph{Source operator.}
\InterfaceTag Fix a separable Hilbert space $\mathcal{H}$ and a trace-class operator $F\in\mathcal{L}^1(\mathcal{H})$.
We refer to $F$ as a \emph{source operator} in the following limited sense: it is the abstract generator whose traces/determinants can be used as a bookkeeping device for cycle statistics and mode factors.

\paragraph{Readout kernels and resolvent traces.}
\InterfaceTag Fix a (declared) family of controllable readout kernels $\mathcal{K}\subset \mathcal{B}(\mathcal{H})$.
For each $K\in\mathcal{K}$, consider the resolvent-trace generating function
\[
T_K(r):=\Tr\!\bigl(K(I-rF)^{-1}\bigr),
\]
whenever it is well-defined, as recorded in Appendix~\ref{app:operator_mother_space}.

\paragraph{Determinant bookkeeping.}
\InterfaceTag The corresponding Fredholm determinant
\[
D(r):=\det(I-rF)
\]
packages orbit/loop trace data through its logarithmic expansion (Appendix~\ref{app:operator_mother_space}).
This provides a single language in which ``prime-cycle'' products and ``spectral-mode'' denominators can be discussed without mixing layers.

\subsection{A finite-dimensional protocol RG instance (kernel view closure)}
\label{subsec:operator_mother_space_dictionary_protocol_rg_instance}

\paragraph{Why a finite-dimensional instance is enough here.}
\InterfaceTag The operator mother space is stated in a trace-class setting for maximal reuse.
However, the protocol RG kernel in the kernel view admits a completely explicit finite-dimensional realization on a fixed $16$-dimensional coarse screen (Section~\ref{subsec:kernel_operator_closure}).
In that setting, all determinant/resolvent objects reduce to ordinary matrix operations and are therefore fully auditable.

\paragraph{Source operator as a protocol RG step.}
\InterfaceTag Fix $n\ge 3$ and let $F$ denote either the unweighted protocol RG operator $F_n$ or the folding-weighted family $\widehat F_n(t)$ acting on $\mathcal{H}:=\RR^{16}$ (Definitions~\ref{def:protocol_rg_operator_Fn} and \ref{def:weighted_protocol_rg_operator}).
Then $F$ is a legitimate source operator in the present dictionary sense: it generates a resolvent $(I-rF)^{-1}$ and a determinant $\det(I-rF)$ that package iterated coarse-grained propagation along the balanced chain.

\paragraph{Readout kernels for one-point and two-point statistics.}
\InterfaceTag A one-point readout (e.g.\ block mean) is a choice of a declared kernel $K^{(1)}$ on $\mathcal{H}$, so that $\Tr(K^{(1)}(I-rF)^{-1})$ packages the corresponding iterated statistics.
For a two-point readout (e.g.\ block variance), the correct operator object lives on the tensor mother space $\mathcal{H}\otimes\mathcal{H}$, with source operator $F\otimes F$ and a declared two-point kernel $K^{(2)}$.
This is the minimal way to keep variance-like observables within the same resolvent-trace language without mixing layers; see Theorem~\ref{thm:block_readout_one_two_point} and Theorem~\ref{thm:resolvent_trace_closure_protocol_rg}.

\subsection{Barrier and rigidity language in the Abel coordinate}
\label{subsec:operator_mother_space_dictionary_barrier}

\paragraph{Abel-first coordinate and interior pole obstruction.}
\InterfaceTag In the Abel-first viewpoint of this paper, one studies generating functions on $|r|<1$ and approaches the boundary along $r\uparrow 1$ (Appendix~\ref{app:protocol_primitives}).
In the operator mother space language, the basic rigidity obstruction is an \emph{interior pole} $|r_0|<1$ in the relevant resolvent-mode factors (Appendix~\ref{app:operator_mother_space}).
\AuditTag This is the operator-theoretic restatement of the pole-barrier discipline (Appendix~\ref{app:trace_pole_barrier_template}) and is used here only as a dictionary for the narrative phrase ``crossing a barrier''.
\AuditTag The canonical Abel finite-part scheme and its scheme-stability within admissible kernel classes is recorded in Appendix~\ref{app:abel_finite_part_notes}.
\AuditTag When ``rates'' or ``responses'' are mentioned in this paper, they should be read as finite differences along the executed update order (or along a declared scan step), not as continuum derivatives unless an explicit matching-layer dictionary is declared.
A reusable finite calculus template for scan-orbit differences, orbit traces, and Abel finite parts is recorded in Appendix~\ref{app:cyclic_calculus_theta_fp}.

\subsection{Finite-family updates: CAP, wormhole pointers, and observer language (dictionary)}
\label{subsec:operator_mother_space_dictionary_updates}

\paragraph{Finite-family updates as operator candidates.}
\InterfaceTag A finite candidate family may be represented as a finite operator family
\[
F_f:=F+\Delta_f,
\qquad f\in\mathcal{F},
\]
where each $\Delta_f$ is finite rank (hence trace class), so that the bookkeeping objects $T_K(r)$ and $D(r)$ remain meaningful within the same mother space (Appendix~\ref{app:operator_mother_space}).
CAP then selects a unique minimizer $f_\star$ of an explicit Motive functional $J:\mathcal{F}\to\RR$ (Appendix~\ref{app:cap_audit_template} and Appendix~\ref{app:wish_motive_definitions}).

\paragraph{Pointer jumps as finite-rank updates.}
\InterfaceTag Protocol-level wormhole-like channels are modeled as pointer jumps (Definition~\ref{def:wormhole_pointer_jump}).
In the operator dictionary, such an added shortcut is encoded as a finite-rank update $\Delta$ (Appendix~\ref{app:operator_mother_space}).
This identification is dictionary-only.

\begin{definition}[Observer (kernel choice; dictionary)]
\label{def:observer_kernel_dictionary}
\InterfaceTag An \emph{observer} is modeled by a choice of readout kernel $K\in\mathcal{K}$ (or, more generally, by a declared kernel family $\mathcal{K}$) used to form the resolvent trace $T_K(r)$.
\AuditTag This definition records a protocol-level fact: changing the readout kernel changes which trace statistics are sampled.
It is not used as a theorem-level premise.
\AuditTag In the mainline contract of this paper, kernel families are treated as explicit finite candidates and are closed by CAP when a quantitative scalar readout is reported (Appendix~\ref{subsec:tick_cap_kernels} and Appendix~\ref{app:kernel_family_cap_closure}).
\end{definition}

\paragraph{Budget-triggered horizons as kernel-limited readout (pointer).}
\InterfaceTag The budget-triggered $\chi$-horizon construction in Section~\ref{subsec:chi_budget_horizon_area_law} can be read in this dictionary as follows:
the reconstructed field $\chi(x)$ is a kernel-dependent readout summary, while the threshold rule selects a region whose residual microstate multiplicity exceeds a declared observer budget (Appendix~\ref{app:protocol_horizon_tick_trap}).
From the mother-space viewpoint, this is a controlled way to say that, for the chosen kernel family, stable traces sample only a coarse boundary summary, while finer interior distinctions correspond to variations within an audited finite candidate family $F_f=F+\Delta_f$ that remain invisible to the observer at fixed budget.

\begin{definition}[Consciousness event (finite-rank update; dictionary)]
\label{def:consciousness_finite_rank_update_dictionary}
\InterfaceTag A \emph{consciousness event} is modeled as a declared finite-rank update $\Delta$ of the source operator, i.e.\ an update step $F\mapsto F+\Delta$ within an audited finite candidate family.
\AuditTag This definition is a controlled way to phrase ``nonlocal shortcuts'' (pointer-jump channels) and ``trajectory deviations'' without importing additional ontic assumptions.
It is not used as a theorem-level premise.
\end{definition}

